\documentclass[11pt]{article}

% ============================================================
%  Research Technical Note
%  Topic: time_block_forecast -- forecasting + scoring spec
%  Companion to tex/ar2_rethink/ar2_rethink.tex.
%  Specifies the computations in run-settings.R (block fit +
%  forecast) and scores.R (lead-time scoring).
% ============================================================

% ---- Packages ----
\usepackage{amsmath, amssymb, amsthm}
\usepackage{mathtools}
\usepackage{bm}
\usepackage[margin=1in]{geometry}
\usepackage{hyperref}
\usepackage{enumitem}
\usepackage{xcolor}

\hypersetup{
  colorlinks=true,
  linkcolor=blue!50!black,
  citecolor=blue!50!black,
  urlcolor=blue!50!black
}

% ---- Notation shortcuts (consistent with ar2_rethink.tex) ----
\newcommand{\R}{\mathbb{R}}
\newcommand{\E}{\mathbb{E}}
\newcommand{\Var}{\operatorname{Var}}
\newcommand{\Cov}{\operatorname{Cov}}
\renewcommand{\Pr}{\mathbb{P}}
\newcommand{\Norm}{\mathcal{N}}
\newcommand{\indicator}{\mathbf{1}}
\newcommand{\given}{\,\vert\,}
\newcommand{\spradius}{\rho_{\!*}}

% ---- Title block ----
\title{
  Time-Block Forecasting and Scoring:\\
  An Implementation Specification\\[4pt]
  \large Research Technical Note: \texttt{time\_block\_forecast}
}
\author{Spatial Skew-$t$ Simulation Study}
\date{2026-05-22}

\begin{document}
\maketitle

% ============================================================
\section{Scope and notation}

\begin{itemize}[leftmargin=1.4em]
  \item This note specifies, equation by equation, the computations
        in \texttt{run-settings.R} (block fit and forecast) and
        \texttt{scores.R} (lead-time scoring) of the
        \texttt{time\_block\_forecast} study. Notation follows the
        companion note \texttt{ar2\_rethink}.
  \item Sites are $s\in\mathcal S$ with $n_s=|\mathcal S|$, all
        observed; calendar time is $t\in\{1,\dots,n_t\}$; latent knots
        are $k\in\{1,\dots,K\}$.
  \item With $K=1$ the Voronoi membership $g(s,t)\equiv 1$ is constant,
        so a single knot serves every site and the knot index is
        suppressed where harmless.
  \item Posterior draws are indexed $m\in\{1,\dots,M\}$, blocks
        $b\in\{1,\dots,B\}$, replicate datasets $d\in\{1,\dots,D\}$.
  \item Write $X_t$ generically for a latent \emph{Gaussian} process
        ($\tau^\star$ or $z^\star$); the observable $(\tau,z)$ follow
        from $X$ by the copula transforms of
        \eqref{eq:inv-tau}--\eqref{eq:inv-z}.
\end{itemize}

% ============================================================
\section{Block geometry and the train--test split}

\begin{itemize}[leftmargin=1.4em]
  \item The record carries $B$ seam times
        $T_o^{(1)}<\cdots<T_o^{(B)}$ and one fixed horizon $H$; block
        $b$ partitions time into a contiguous train window and a
        forecast window
        \begin{equation}
          \mathcal T_b^{\mathrm{tr}}
            = \{1,\dots,T_o^{(b)}\},
          \qquad
          \mathcal T_b^{\mathrm{te}}
            = \big\{\,T_o^{(b)}+h \;:\; h=1,\dots,H \,\big\}.
          \label{eq:windows}
        \end{equation}
  \item The lead time $h$ and the calendar time are linked by
        $t = T_o^{(b)}+h$; the study fixes
        \begin{equation}
          n_t = 200,\quad B = 5,\quad
          \big(T_o^{(b)}\big)_{b=1}^{5} = (50,80,110,140,170),\quad
          H = 15 .
          \label{eq:constants}
        \end{equation}
  \item Every site is observed, so the holdout is purely temporal: the
        training data of block $b$ is
        $\{y(s,t):s\in\mathcal S,\ t\in\mathcal T_b^{\mathrm{tr}}\}$ and
        the target is
        $\{y(s,t):s\in\mathcal S,\ t\in\mathcal T_b^{\mathrm{te}}\}$.
  \item Adjacent seams are spaced by
        $T_o^{(b+1)}-T_o^{(b)} = 30 > H$, so no forecast window reaches
        the next seam and the $B$ per-block scores are treated as
        near-independent in \eqref{eq:leadcurve}.
\end{itemize}

% ============================================================
\section{Posterior fit on the prefix}

\begin{itemize}[leftmargin=1.4em]
  \item For each $(b,d,\text{method})$ the backend draws $M$ samples
        from the prefix posterior
        \begin{equation}
          \pi_b\big(\theta,\,\{X^\star_{k,t}\}_{t\le T_o^{(b)}}\big)
          = \pi\big(\theta,\,\{X^\star_{k,t}\}_{t\le T_o^{(b)}}
              \,\big\vert\,
              \{y(s,t):s\in\mathcal S,\ t\le T_o^{(b)}\}\big).
          \label{eq:posterior}
        \end{equation}
  \item Draw $m$ supplies the parameter bundle
        \begin{equation}
          \theta^{(m)}
          = \big(\beta^{(m)},\,\lambda^{(m)},\,
                 \alpha^{(m)},\,\beta_\tau^{(m)},\,
                 \rho^{(m)},\,\nu^{(m)},\,\gamma^{(m)},\,
                 \phi_\tau^{(m)},\,\phi_z^{(m)}\big)
          \label{eq:thetabundle}
        \end{equation}
        together with the imputed observable latent trajectories
        $\tau^{(m)}_{k,t},\,z^{(m)}_{k,t}$ for $t\le T_o^{(b)}$.
  \item The gamma copula uses the halved hyperparameters
        $a^{(m)} = \alpha^{(m)}/2$ and $b^{(m)} = \beta_\tau^{(m)}/2$,
        matching the generator's reparameterization.
  \item The i.i.d.\ method fixes $\phi_\tau=\phi_z=\bm 0$ and disables
        temporal pooling; the AR(2) method estimates
        $\phi_\tau^{(m)},\phi_z^{(m)}\in\R^2$.
\end{itemize}

% ============================================================
\section{Forecasting the held-out window (\texttt{forecast\_block})}

\subsection{Seam-state extraction}

\begin{itemize}[leftmargin=1.4em]
  \item The recursion runs on the Gaussian scale, so the imputed
        \emph{observable} seam values at $t\in\{T_o-1,\,T_o\}$ are
        mapped back by the forward copula transforms
        \begin{align}
          \tau^{\star(m)}_{k,t}
            &= \Phi^{-1}\!\big(G_{a^{(m)},b^{(m)}}(\tau^{(m)}_{k,t})\big),
          \label{eq:fwd-tau}\\
          z^{\star(m)}_{k,t}
            &= \Phi^{-1}\!\big(H_{\sigma^{(m)}_{z,t}}(z^{(m)}_{k,t})\big),
          \qquad
          \sigma^{(m)}_{z,t} = \big(\tau^{(m)}_{k,t}\big)^{-1/2},
          \label{eq:fwd-z}
        \end{align}
        with $\Phi$ the standard normal CDF, $G_{a,b}$ the gamma CDF,
        and $H_\sigma$ the half-normal CDF of scale $\sigma$.
  \item The seam state $(X^{\star(m)}_{k,T_o-1},X^{\star(m)}_{k,T_o})$
        is taken from draw $m$'s own imputed trajectory, never from the
        stationary law $\Norm(\bm 0,\Sigma)$; by Corollary~1 of the
        companion note that choice is what preserves the AR(2) signal.
\end{itemize}

\subsection{Latent AR(2) recursion}

\begin{itemize}[leftmargin=1.4em]
  \item For draw $m$, knot $k$, and process $X\in\{\tau^\star,z^\star\}$
        with its own coefficients
        $\phi^{(m)}=(\phi^{(m)}_1,\phi^{(m)}_2)$, the window is recursed
        for $h=1,\dots,H$:
        \begin{equation}
          X^{\star(m)}_{k,T_o+h}
          = \phi^{(m)}_1 X^{\star(m)}_{k,T_o+h-1}
          + \phi^{(m)}_2 X^{\star(m)}_{k,T_o+h-2}
          + \sigma^{(m)}\,\xi^{(m)}_{k,h},
          \qquad
          \xi^{(m)}_{k,h}\overset{\text{iid}}{\sim}\Norm(0,1).
          \label{eq:recurse}
        \end{equation}
  \item The innovation standard deviation is fixed by Yule--Walker at
        unit marginal variance ($\gamma_0=1$), so the latent variance
        is preserved across the forecast:
        \begin{equation}
          \gamma_1 = \frac{\phi^{(m)}_1}{1-\phi^{(m)}_2},
          \qquad
          \gamma_2 = \phi^{(m)}_1\gamma_1 + \phi^{(m)}_2,
          \qquad
          \big(\sigma^{(m)}\big)^2
            = 1 - \phi^{(m)}_1\gamma_1 - \phi^{(m)}_2\gamma_2 .
          \label{eq:yw}
        \end{equation}
  \item A draw whose $\phi^{(m)}$ leaves the stationarity triangle
        $\{|\phi_2|<1,\ \phi_1+\phi_2<1,\ \phi_2-\phi_1<1\}$ is clamped
        to $\sigma^{(m)}=1$, capping the non-stationary recursion at the
        i.i.d.\ marginal.
\end{itemize}

\subsection{The i.i.d.\ baseline}

\begin{itemize}[leftmargin=1.4em]
  \item The i.i.d.\ method discards the seam state and draws each
        latent forecast slice fresh from the stationary marginal:
        \begin{equation}
          \tau^{\star(m)}_{k,T_o+h},\;
          z^{\star(m)}_{k,T_o+h}
          \overset{\text{iid}}{\sim} \Norm(0,1),
          \qquad h = 1,\dots,H .
          \label{eq:iid}
        \end{equation}
  \item Equation~\eqref{eq:iid} is exactly the $h\to\infty$ limit of
        \eqref{eq:recurse}, so the two methods differ \emph{only}
        through the seam-state conditioning.
\end{itemize}

\subsection{Copula transforms and the spatial draw}

\begin{itemize}[leftmargin=1.4em]
  \item After the recursion the latent window is mapped to the
        observable scale by the inverse copula transforms
        \begin{align}
          \tau^{(m)}_{k,T_o+h}
            &= G^{-1}_{a^{(m)},b^{(m)}}\!\big(
                 \Phi(\tau^{\star(m)}_{k,T_o+h})\big),
          \label{eq:inv-tau}\\
          z^{(m)}_{k,T_o+h}
            &= H^{-1}_{\sigma^{(m)}_z}\!\big(
                 \Phi(z^{\star(m)}_{k,T_o+h})\big),
          \qquad
          \sigma^{(m)}_z = \big(\tau^{(m)}_{k,T_o+h}\big)^{-1/2},
          \label{eq:inv-z}
        \end{align}
        which invert \eqref{eq:fwd-tau}--\eqref{eq:fwd-z} and apply
        \emph{after} the recursion.
  \item The $\tau$--$z$ coupling is enforced within the draw:
        $\sigma^{(m)}_z$ in \eqref{eq:inv-z} uses the same draw's
        forecasted $\tau^{(m)}_{k,T_o+h}$, not an independent marginal.
  \item For each lead $h$ (time $t=T_o+h$) the spatial field is drawn
        from the skew-$t$ observation model with $K=1$ membership
        $g\equiv 1$:
        \begin{equation}
          \hat y^{(m)}(s,t)
          = x(s,t)^\top\beta^{(m)}
          + \lambda^{(m)} z^{(m)}_{1,t}
          + \varepsilon^{(m)}(s,t),
          \label{eq:obs-forecast}
        \end{equation}
        \begin{equation}
          \varepsilon^{(m)}(\cdot,t)\sim
            \Norm\!\Big(\bm 0,\ \big(\tau^{(m)}_{1,t}\big)^{-1} C^{(m)}\Big),
          \qquad
          C^{(m)}_{ss'} = C\big(s,s';\rho^{(m)},\nu^{(m)},\gamma^{(m)}\big).
          \label{eq:spatial}
        \end{equation}
  \item Operationally \eqref{eq:spatial} is realized as
        $\varepsilon^{(m)}(\cdot,t) = L^{(m)}\zeta$, with
        $L^{(m)}L^{(m)\top}=C^{(m)}$ the Cholesky factor and
        $\zeta\sim\Norm\!\big(\bm0,(\tau^{(m)}_{1,t})^{-1}I\big)$.
  \item The block output is the predictive sample array
        $\big\{\hat y^{(m)}(s,T_o^{(b)}+h)\big\}$ of shape
        $M\times n_s\times H$, paired with the validation truth
        $y^{\mathrm{val}}_b(s,h) = y(s,T_o^{(b)}+h)$.
  \item Averaging over $m=1,\dots,M$ integrates \eqref{eq:thetabundle}
        and the seam state jointly, so the predictive carries both
        parameter and state uncertainty with no plug-in of posterior
        means.
\end{itemize}

% ============================================================
\section{Score computation (\texttt{scores.R})}

\begin{itemize}[leftmargin=1.4em]
  \item Each $(b,s,h)$ predictive is the $M$-point empirical measure
        \begin{equation}
          \hat F_{b,s,h}
          = \frac1M \sum_{m=1}^M
            \delta_{\,\hat y^{(m)}(s,\,T_o^{(b)}+h)} ,
          \label{eq:empmeasure}
        \end{equation}
        and all scores below are functionals of $\hat F_{b,s,h}$ and the
        truth $y^{\mathrm{val}}_b(s,h)$.
\end{itemize}

\subsection{Lead-time CRPS}

\begin{itemize}[leftmargin=1.4em]
  \item The per-site CRPS uses the empirical-distribution estimator
        \begin{equation}
          \mathrm{CRPS}\big(\hat F_{b,s,h},\,y\big)
          = \frac1M\sum_{m=1}^M \big|\hat y^{(m)}-y\big|
          - \frac1{2M^2}\sum_{m=1}^M\sum_{m'=1}^M
            \big|\hat y^{(m)}-\hat y^{(m')}\big|,
          \label{eq:crps}
        \end{equation}
        with $\hat y^{(m)}=\hat y^{(m)}(s,T_o^{(b)}+h)$; this is the
        $d=1$ case of the energy score.
  \item The block-$b$, lead-$h$ score is the site average
        \begin{equation}
          S^{\mathrm{CRPS}}(h;b,d)
          = \frac1{n_s}\sum_{s\in\mathcal S}
            \mathrm{CRPS}\big(\hat F_{b,s,h},\,
              y^{\mathrm{val}}_b(s,h)\big).
          \label{eq:crps-block}
        \end{equation}
\end{itemize}

\subsection{Lead-time Brier}

\begin{itemize}[leftmargin=1.4em]
  \item Thresholds are the empirical quantiles of the block's
        validation field at the probability grid
        $\bm p = (0.90,0.92,0.94,0.95,0.96,0.98,0.99)$:
        \begin{equation}
          c_{b,q}
          = \hat Q_{p_q}\!\big(\{\,y(s,t):
              s\in\mathcal S,\ t\in\mathcal T_b^{\mathrm{te}}\,\}\big),
          \qquad q = 1,\dots,|\bm p| .
          \label{eq:thresholds}
        \end{equation}
  \item The forecast exceedance probability and the lead-$h$,
        threshold-$q$ Brier score are
        \begin{align}
          \hat p_b(s,h;q)
            &= \frac1M\sum_{m=1}^M
               \indicator\big\{\hat y^{(m)}(s,T_o^{(b)}+h) > c_{b,q}\big\},
          \label{eq:phat}\\
          S^{\mathrm{BS}}(h,q;b,d)
            &= \frac1{n_s}\sum_{s\in\mathcal S}
               \Big(\indicator\big\{y^{\mathrm{val}}_b(s,h)>c_{b,q}\big\}
                    - \hat p_b(s,h;q)\Big)^2 .
          \label{eq:brier}
        \end{align}
\end{itemize}

\subsection{Joint energy and variogram scores}

\begin{itemize}[leftmargin=1.4em]
  \item For site $s$ the held-out window is the length-$H$ vector
        $\bm y_{b,s} = \big(y(s,T_o^{(b)}+1),\dots,y(s,T_o^{(b)}+H)\big)
        \in\R^H$, with predictive sample
        $\hat{\bm y}^{(m)}_{b,s}\in\R^H$.
  \item The energy score (sample estimator, $\|\cdot\|$ the Euclidean
        norm on $\R^H$) is
        \begin{equation}
          \widehat{\mathrm{ES}}(b,s)
          = \frac1M\sum_{m=1}^M
              \big\|\hat{\bm y}^{(m)}_{b,s} - \bm y_{b,s}\big\|
          - \frac1{2M^2}\sum_{m=1}^M\sum_{m'=1}^M
              \big\|\hat{\bm y}^{(m)}_{b,s}
                    - \hat{\bm y}^{(m')}_{b,s}\big\| .
          \label{eq:es}
        \end{equation}
  \item The variogram score of order $p=\tfrac12$, unit weights, and
        $i,j$ indexing lead time is
        \begin{equation}
          \widehat{\mathrm{VS}}(b,s)
          = \sum_{i=1}^H\sum_{j=1}^H
            \Big(\,
              |y_{b,s,i}-y_{b,s,j}|^{p}
              - \frac1M\sum_{m=1}^M
                |\hat y^{(m)}_{b,s,i}-\hat y^{(m)}_{b,s,j}|^{p}
            \,\Big)^2 .
          \label{eq:vs}
        \end{equation}
  \item Each is reduced to one number per block by site averaging,
        \begin{equation}
          S^{\mathrm{ES}}(b,d) = \frac1{n_s}\sum_{s}\widehat{\mathrm{ES}}(b,s),
          \qquad
          S^{\mathrm{VS}}(b,d) = \frac1{n_s}\sum_{s}\widehat{\mathrm{VS}}(b,s).
          \label{eq:joint-block}
        \end{equation}
  \item In \eqref{eq:es}--\eqref{eq:vs} the lead time supplies the
        multivariate dimension, so ES and VS credit the temporal
        dependence that the pointwise CRPS of \eqref{eq:crps} cannot
        resolve.
\end{itemize}

\subsection{Lead-time curve}

\begin{itemize}[leftmargin=1.4em]
  \item The Brier lead curve first averages \eqref{eq:brier} over the
        threshold grid,
        \begin{equation}
          S^{\mathrm{BS}}(h;b,d)
          = \frac1{|\bm p|}\sum_{q=1}^{|\bm p|} S^{\mathrm{BS}}(h,q;b,d).
          \label{eq:brier-leadavg}
        \end{equation}
  \item For a univariate score
        $S\in\{S^{\mathrm{CRPS}},S^{\mathrm{BS}}\}$ the per-block mean
        over the $D$ replicate datasets is
        \begin{equation}
          \bar S_b(h) = \frac1D\sum_{d=1}^D S(h;b,d).
          \label{eq:blockmean}
        \end{equation}
  \item The lead-time curve and its standard error follow from the
        dispersion of the $B$ per-block means,
        \begin{equation}
          \bar S(h) = \frac1B\sum_{b=1}^B \bar S_b(h),
          \qquad
          \mathrm{SE}(h)
          = \frac1{\sqrt B}
            \Big(\frac1{B-1}\sum_{b=1}^B
                 \big(\bar S_b(h)-\bar S(h)\big)^2\Big)^{1/2}.
          \label{eq:leadcurve}
        \end{equation}
  \item Because the seam spacing $30>H$ makes the $B$ summands of
        \eqref{eq:leadcurve} near-independent, $\mathrm{SE}(h)$ is a
        valid uncertainty band for $\bar S(h)$.
\end{itemize}

% ============================================================
\section{Verification conditions}

\begin{itemize}[leftmargin=1.4em]
  \item \textbf{$H=0$ parity.} Re-scoring time $T_o$ must give the
        AR(2) and i.i.d.\ methods the same predictive; any difference
        is a seam-indexing error in \eqref{eq:fwd-tau}--\eqref{eq:fwd-z}.
  \item \textbf{$\phi=\bm 0$ collapse.} A draw with
        $\phi^{(m)}=\bm 0$ must reduce \eqref{eq:recurse} to
        \eqref{eq:iid} immediately; residual autocorrelation signals a
        double-counted seam state.
  \item \textbf{Long-horizon limit.} With stationary $\phi^{(m)}$ and
        $H\gg h^\star$, the forecasted latent slices must reproduce the
        $\Norm(0,1)$ marginal, confirming the Yule--Walker
        normalization \eqref{eq:yw}.
\end{itemize}

\end{document}
